\chapter{Modelo Relacional}

\section[Conceptos fundamentales]{\emj{🧠}~Conceptos fundamentales}

El modelo relacional, propuesto por E.F.~Codd en 1970, representa datos como tablas (relaciones matemáticas).
Cada fila es una tupla única; cada columna es un atributo con dominio definido.
Las relaciones entre tablas se expresan mediante llaves, no mediante punteros físicos.

\begin{teoriabox}{Terminología formal vs coloquial}
\begin{itemize}
  \item \textbf{Relación} (formal) = \textbf{tabla} (coloquial)
  \item \textbf{Tupla} (formal) = \textbf{fila / registro}
  \item \textbf{Atributo} (formal) = \textbf{columna / campo} con tipo de dato fijo
  \item \textbf{Dominio:} conjunto de valores válidos para un atributo (ej.~enteros positivos)
  \item \textbf{Grado:} número de atributos (columnas) de la relación
  \item \textbf{Cardinalidad:} número de tuplas (filas) actuales en la relación
\end{itemize}
La fortaleza del modelo relacional es su independencia física: el motor puede cambiar cómo almacena los datos sin afectar las consultas SQL.
\end{teoriabox}

\section[Llaves]{\emj{🧠}~Llaves}

Las llaves son el mecanismo de identificación y relación entre tablas.
Una tabla bien diseñada siempre tiene una PK definida; las FK crean el grafo de relaciones del esquema.
Elegir entre surrogate y natural key afecta rendimiento, legibilidad y mantenimiento a largo plazo.

\begin{itemize}
  \item \textbf{Superkey:} conjunto de atributos que identifica tuplas únicamente
  \item \textbf{Candidate key:} superkey mínima (sin atributos redundantes)
  \item \textbf{Primary Key (PK):} candidate key elegida; NOT NULL + UNIQUE garantizados
  \item \textbf{Foreign Key (FK):} atributo que referencia la PK de otra tabla
  \item \textbf{Surrogate key:} PK artificial (SERIAL/UUID) sin significado de negocio
  \item \textbf{Natural key:} PK con significado de negocio (RFC, email, CURP)
\end{itemize}

\begin{tipbox}{Surrogate vs Natural Key}
Prefiere surrogate keys (SERIAL/UUID) para tablas con relaciones: son inmutables, sin significado de negocio que cambie y más eficientes en joins.
Usa natural keys solo cuando el atributo es \emph{verdaderamente} estable y único (ej. ISBN de un libro, código IATA de un aeropuerto).
\end{tipbox}

\section[Integridad referencial]{\emj{🧠}~Integridad referencial}

La integridad referencial garantiza que una FK siempre apunte a un registro que existe en la tabla padre, o sea NULL si está permitido.
Sin FK enforced, es posible tener huérfanos: órdenes sin cliente, facturas sin producto.
El DBMS rechaza operaciones que violen la restricción.

Las acciones \texttt{ON DELETE} / \texttt{ON UPDATE} definen qué pasa con los hijos cuando cambia el padre:
\begin{itemize}
  \item \textbf{CASCADE:} propaga el cambio a los hijos automáticamente
  \item \textbf{SET NULL:} pone NULL en la FK del hijo (la columna debe admitir NULL)
  \item \textbf{RESTRICT / NO ACTION:} rechaza la operación si existen hijos (default)
\end{itemize}

\begin{lstlisting}[caption={PK, FK e integridad referencial}]
CREATE TABLE departamentos (
    id     SERIAL PRIMARY KEY,
    nombre VARCHAR(80) NOT NULL UNIQUE
);

CREATE TABLE empleados (
    id       SERIAL PRIMARY KEY,
    nombre   VARCHAR(100) NOT NULL,
    id_depto INTEGER REFERENCES departamentos(id)
             ON DELETE SET NULL
             ON UPDATE CASCADE
);

-- FK viola integridad si el depto no existe:
INSERT INTO empleados (nombre, id_depto) VALUES ('Carlos', 999);
-- ERROR: insert or update on table "empleados" violates foreign key constraint
\end{lstlisting}

\begin{lstlisting}[caption={Cardinalidades 1:1, 1:N, M:N}]
-- 1:N — un depto tiene muchos empleados (FK en "muchos")
-- (ya modelado arriba con id_depto en empleados)

-- M:N — un empleado puede estar en muchos proyectos
CREATE TABLE proyectos (
    id     SERIAL PRIMARY KEY,
    nombre VARCHAR(100) NOT NULL
);

CREATE TABLE empleados_proyectos (
    id_empleado INTEGER REFERENCES empleados(id)   ON DELETE CASCADE,
    id_proyecto INTEGER REFERENCES proyectos(id)   ON DELETE CASCADE,
    rol         VARCHAR(50),
    PRIMARY KEY (id_empleado, id_proyecto)
);
\end{lstlisting}

\begin{entrevistabox}
\begin{enumerate}
  \item ¿Cuál es la diferencia entre Primary Key y Unique Key?
  \item ¿Qué sucede con los registros hijos cuando se borra el padre con \texttt{ON DELETE CASCADE}?
  \item ¿Cuándo preferirías una surrogate key sobre una natural key?
\end{enumerate}
\end{entrevistabox}
